\documentclass[11pt,a4paper]{article}

\usepackage[utf8]{inputenc}
\usepackage[T1]{fontenc}
\usepackage{lmodern}
\usepackage{amsmath,amssymb,amsthm}
\usepackage{hyperref}
\usepackage{geometry}
\usepackage{doi}
\geometry{margin=2.5cm}

\input{glyphtounicode}
\pdfgentounicode=1

\newtheorem{theorem}{Theorem}[section]
\newtheorem{lemma}[theorem]{Lemma}
\newtheorem{corollary}[theorem]{Corollary}
\newtheorem{proposition}[theorem]{Proposition}
\newtheorem{definition}[theorem]{Definition}
\newtheorem{remark}[theorem]{Remark}

\newcommand{\bstar}{\beta^*}
\newcommand{\dpair}{\delta_{\mathrm{pair}}}
\newcommand{\Ntrl}{\mathcal{N}_{\mathrm{trl}}}
\newcommand{\Aadm}{\mathcal{A}}
\newcommand{\Obs}{\mathcal{O}}
\newcommand{\GBI}{G_{\mathrm{BI}}}
\newcommand{\SBI}{S_{\mathrm{BI}}}
\newcommand{\ImH}{\operatorname{Im}\mathbb{H}}
\newcommand{\cchi}{c_{\mathrm{BI}}}

\begin{document}

  \title{Observable-Rank Stability under Vertical Non-Injectivity:\\
  Closing the Fibre-Structure Conditionality in the
  Born--Infeld-to-Cascade Chain\\[6pt]
    {\large O24 --- Spectral Admissibility Sub-programme}}

  \author{J\'er\^ome Beau\\
  \small Independent Researcher\\
  \small \href{mailto:jerome.beau@cosmochrony.org}{jerome.beau@cosmochrony.org}}

  \date{2026\\Version 1.3}

  \maketitle

  \begin{abstract}
    O18 establishes Born--Infeld parity as a covariance of the response family of $\SBI[\chi]$ and states
    the identification of the parity orbit with a fibre of the non-injective projection $\Pi$
    as an open problem; the hypothesis that parity is the only global symmetry of $\SBI$ is a
    clause of that problem, not a theorem.
    O22 proved the shell-alignment conjecture of O21.
    O23 proves that the traceless sector of a supplied spinor carrier
    $V_\rho \cong \mathbb{C}^2$ is $\mathfrak{su}(2)$, of real dimension exactly~$3$; the
    carrier selection and the identification of $\Sigma_c$ with that dimension remain open, so
    the threshold $\Sigma_c(n_3) = 3$ is a supplied selection rule.
    Together, these results locate the integer~$3$ in $\Sigma_c(n_3) = 3$ as a property of
    $\operatorname{Im}\Pi \cap \Ntrl$, not of the cardinality of the fibres.
    The programme then requires only that
    any symmetry of $\SBI$ acts \emph{vertically} with respect to $\Pi$, i.e.\ does not enlarge
    the real rank of the admissible neutral traceless sector.
    The present paper proves this \emph{verticality condition} from Born--Infeld
    admissibility and the rank-three carrier supplied by~O23, without requiring a
    classification of $\mathcal{C}_{\mathrm{eff}}$.
    As a corollary, conditional on the supplied carrier and fibre hypotheses, the
    fibre-structure conditionality in the $\cchi \to \dpair$ segment reduces to those supplied
    inputs: the observable rank is independent of fibre cardinality.
    This does not establish the separate
    $\dpair \to \bstar$ capacity-to-rate step, for which no native
    Heisenberg growth carrier is defined in the corpus.
    Non-injectivity may increase microscopic multiplicity, but cannot enlarge the admissible
    observable structure.
  \end{abstract}

  \medskip
  \noindent\textbf{Keywords.}
  Cosmochrony; spectral admissibility; non-injective projection; verticality condition;
  observable rank stability; spinor carrier; Born--Infeld constraint;
  fibre structure; cascade exponent~$\bstar$

  \clearpage
  \tableofcontents

  \clearpage
  \section{Position of the Problem after O22--O23}
    \label{sec:position}

    \subsection{What O18 leaves open and what the mechanism requires}
      \label{sec:o18}

      O18~\cite{Beau2026a22} establishes that the Born--Infeld action $\SBI[\chi]$ is even
      (O18 Lemma~2.1) and that parity acting simultaneously on configuration and probe is an
      covariance of the Born--Infeld response family (even-order responses preserved, odd-order reversed) (O18 Proposition~2.6).
      Evenness does not force fixed-probe indiscernibility of $\chi$ and $-\chi$
      (O18 Proposition~2.4), so the parity orbit $\{\chi, -\chi\}$ is not forced into the
      fibres of $\Pi$: the fibre identification is a typed classification problem
      (O18 Problem~2.8), and the assumption that parity is the only global effective symmetry
      of $\SBI$ is a clause of that open problem (O18 Remark~2.9), not a theorem.

      The condition ``parity is the only symmetry'' would, if established, guarantee minimality
      of a parity fibre $\{\chi, -\chi\}$, ensuring that the non-injectivity of $\Pi$ does not
      create additional observable directions beyond those already present.
      It is, however, stronger than what the present mechanism requires.
      The programme does not require small fibres; it requires that additional fibre structure, if
      present, be \emph{invisible at the level of admissible observables}.

    \subsection{The key distinction: fibre cardinality versus observable rank}
      \label{sec:distinction}

      The results of O22 and O23 clarify the precise invariant that controls the fibre-side
      $\cchi \to \dpair$ segment.
      O23~\cite{Beau2026a27} identifies this invariant as
      \[
        \dim_{\mathbb{R}}\bigl(\operatorname{Im}\Pi \cap \Ntrl\bigr) = 3,
      \]
      where $\Ntrl$ denotes the admissible neutral traceless sector.
      This is a property of the \emph{image} of $\Pi$, not of its fibres.
      The integer~$3$ is the real dimension of $\mathfrak{su}(2) \cong \ImH$, the traceless sector of a
      supplied spinor carrier $V_\rho \cong \mathbb{C}^2$ (O23 Theorem~3.1); the selection of
      that carrier and the identification of $\Sigma_c$ with its adjoint dimension are open
      bridges of O23, so the rank-three input is a supplied hypothesis throughout this paper.

      Two configurations that lie in the same fibre of $\Pi$ are, by definition, mapped to the same
      observable.
      Whether the fibre contains two elements or a continuum, the image of $\Pi$ --- and hence the
      admissible observable structure --- is unchanged.
      The cardinality of the fibres and the rank of the image are therefore independent quantities,
      and it is the rank that governs the mechanism.

    \subsection{The two levels of $\Pi$: kernel and image}
      \label{sec:twolevels}

      The projection $\Pi$ admits a natural decomposition into two structurally distinct levels.

      The \emph{kernel level} $\ker\Pi$ encodes the microscopic multiplicity: the set of substrate
      configurations $\chi$ that are indistinguishable at the level of admissible observables.
      A larger kernel means more microscopic configurations projecting to the same effective state.
      This does not generate new observable directions; it generates projection residuals, i.e.\
      substrate degrees of freedom that are compressed by $\Pi$ and manifest as fluctuations in
      the effective description.

      The \emph{image level} $\operatorname{Im}\Pi$ encodes the effective observable structure: the
      set of states that are distinguishable by at least one admissible observable.
      Its intersection with $\Ntrl$ has real rank~$3$ under the supplied spinor-carrier
      hypothesis of O23 (Theorem~3.1 there), and this rank is what enters the supplied
      selection rule $\Sigma_c(n_3) = 3$.
      Its insertion into a capacity-to-rate law is a separate step.

      The key structural principle is:
      \begin{quote}
        \emph{Fluctuations and observables coexist in the same effective space.
        More residuals do not add new directions to that space; they increase the amplitude of
        projection noise within the directions already present.}
      \end{quote}

    \subsection{Scope of O24}
      \label{sec:scope}

      The present paper proves the \emph{verticality condition}: every symmetry of $\SBI$
      compatible with admissibility acts vertically with respect to $\Pi$, in the sense that it
      does not increase the real rank of $\operatorname{Im}\Pi \cap \Ntrl$.
      The proof uses the rank-three carrier of O23 (Theorem~3.1 there, conditional on the
      supplied spinor carrier) as its central ingredient and
      does not require a global classification of $\mathcal{C}_{\mathrm{eff}}$.

      The paper does not address the remaining open task of the programme: the extended numerical
      campaign at large primes $q \in \{101, 151, 211\}$ across multiple conjugate pairs, required
      to confirm the shell-alignment prediction and reduce the residual distance
      $\mathrm{dist}(n_3^{\mathrm{obs}}, \mathbb{Z}) \approx 0.20$ observed at $q = 61$.

  \section{The Observable Rank as the Invariant of the Chain}
    \label{sec:rank}

    \subsection{Recapitulation: the conditional rank-three carrier of O23}
      \label{sec:o23recap}

      We recall the central result of O23~\cite{Beau2026a27}.

      \begin{theorem}[O23 Theorem~3.1, Three neutral directions from a spinor carrier]
        \label{thm:o23}
        Suppose the admissible neutral sector is carried by an irreducible two-dimensional
        unitary representation $V_\rho \cong \mathbb{C}^2$ (a supplied input).
        Then its traceless sector is $\mathfrak{su}(2) \cong \ImH$, of real dimension
        exactly~$3$.
      \end{theorem}

      Two bridges remain open in O23: nothing derives the carrier $V_\rho \cong \mathbb{C}^2$
      from Born--Infeld parity (carrier selection), and no theorem identifies $\Sigma_c$ with
      $\dim \mathfrak{su}(V_\rho)$ (observable identification).
      The threshold $\Sigma_c(n_3) = 3$ is therefore a supplied selection rule, definitional at
      the level of O21, and every statement of the present paper that consumes the rank-three
      input is conditional on the supplied carrier and on this identification.
      Under the carrier hypothesis the rank cannot exceed~$3$ within $\mathfrak{su}(V_\rho)$,
      which is what the verticality argument of Section~\ref{sec:verticality} uses; the
      conditionality is inherited there.

    \subsection{Independence of the observable rank from the fibre structure}
      \label{sec:independence}

      The following lemma makes the rank--kernel decoupling explicit.

      \begin{lemma}[Rank--kernel decoupling]
        \label{lem:decoupling}
        The real rank of $\operatorname{Im}\Pi \cap \Ntrl$ is independent of the cardinality of
        $\ker\Pi$.
      \end{lemma}

      \begin{proof}
        The rank is a property of the image of $\Pi$ as a linear map; it is determined by the dimension
        of the span of $\Pi(\Aadm)$ within $\Ntrl$.
        Enlarging $\ker\Pi$ --- i.e.\ increasing the number of substrate configurations that project to
        the same observable --- leaves the image $\Pi(\Aadm)$ and its intersection with $\Ntrl$
        unchanged.
        The rank therefore depends only on the structure of $\operatorname{Im}\Pi$, not on the
        structure of $\ker\Pi$.
      \end{proof}

      \begin{remark}
        Lemma~\ref{lem:decoupling} is linear-algebraic and holds independently of any assumption on
        the symmetry group $\GBI$.
        It records the conceptual separation between the two levels of $\Pi$ identified in
        Section~\ref{sec:twolevels}: kernel structure governs microscopic multiplicity and
        fluctuation amplitude; image structure governs the observable mechanism.
      \end{remark}

    \subsection{Residuals and fluctuations as kernel-level phenomena}
      \label{sec:residuals}

      When the fibre of $\Pi$ above an observable $o$ contains more than the minimal orbit
      $\{\chi, -\chi\}$, the additional configurations constitute substrate degrees of freedom that
      are invisible to admissible observables.
      These degrees of freedom are not lost; they appear as projection residuals in the effective
      description.

      Formally, let $F_o = \Pi^{-1}(o)$ be the fibre above $o$.
      The projection decomposes as
      \[
        \chi = \chi_{\parallel} + \chi_{\perp},
      \]
      where $\chi_{\parallel}$ is the component of $\chi$ in the directions visible to $\Pi$
      (i.e.\ transverse to $\ker\Pi$) and $\chi_{\perp} \in \ker\Pi$ is the residual.
      The observable $o = \Pi(\chi)$ depends only on $\chi_{\parallel}$.
      A larger fibre means a larger space of possible $\chi_{\perp}$ values; it does not modify
      the effective dynamics of $\chi_{\parallel}$.

      In particular, the cascade exponent $\dpair$, the saturation rank $n_3$, and the transfer
      $\dpair \to \bstar$ all operate at the level of $\chi_{\parallel}$ --- the observable sector.
      Projection residuals $\chi_{\perp}$ appear as fluctuations around the effective dynamics but
      do not alter the mechanism.
      Fluctuations and observables coexist in the same effective space; additional residuals
      increase the noise amplitude within that space --- its statistical occupation measure ---
      not its dimensionality.
      This is consistent with the identification in O14~\cite{Beau2026a18} Section~6.4 of the
      residual correction $\epsilon(q)$ as the observable imprint of projection residuals, controlled
      by the phase variance $\mathrm{Var}(\theta)(q)$: more residuals shift the statistical weight
      within the observable space, leaving its rank unchanged.

  \section{The Verticality Lemma}
    \label{sec:verticality}

    \subsection{Setup and statement}
      \label{sec:setup}

      Let $\GBI$ denote the group of symmetries of $\SBI[\chi]$ compatible with Born--Infeld
      admissibility.
      A symmetry $g \in \GBI$ is said to be \emph{admissible} if $g(\Aadm) \subset \Aadm$, i.e.\
      it maps the admissible sector to itself.
      All elements of $\GBI$ are admissible by definition; the property is stated explicitly here
      because it is the key hypothesis used in the proof of Lemma~\ref{lem:verticality}.
      A symmetry $g \in \GBI$ is said to act \emph{vertically} with respect to $\Pi$ if it maps
      each fibre $\Pi^{-1}(o)$ to itself, i.e.\ $\Pi(g\chi) = \Pi(\chi)$ for all $\chi \in \Aadm$.
      A symmetry acts \emph{transversally} if there exists $\chi \in \Aadm$ such that
      $\Pi(g\chi) \neq \Pi(\chi)$, and moreover
      $\Pi(g\chi) \notin \mathrm{span}\bigl(\Pi(\Aadm)\bigr)$.

      A transversal symmetry therefore generates a new linearly independent direction in
      $\operatorname{Im}\Pi \cap \Ntrl$: it strictly increases the real rank of
      $\operatorname{Im}\Pi \cap \Ntrl$.

      \begin{lemma}[Verticality of admissible symmetries]
        \label{lem:verticality}
        Let $g \in \GBI$ be a symmetry of $\SBI[\chi]$ compatible with Born--Infeld admissibility.
        Then $g$ acts vertically with respect to $\Pi$: it does not increase the real rank of
        $\operatorname{Im}\Pi \cap \Ntrl$.
      \end{lemma}

      \begin{proof}
        Suppose for contradiction that $g$ acts transversally.
        Then $g$ generates a new linearly independent direction in
        $\operatorname{Im}\Pi \cap \Ntrl$.

        Since $g \in \GBI$ is admissible, $g(\Aadm) = \Aadm$: any vector generated by the action
        of $g$ on $\Aadm$ remains in $\Aadm$, and in particular satisfies the Born--Infeld
        admissibility constraint.
        Such a direction therefore belongs to the admissible neutral traceless sector characterised
        in Theorem~\ref{thm:o23}.

        But under the supplied carrier hypothesis, Theorem~\ref{thm:o23} (O23) bounds the
        admissible neutral traceless sector by $\mathfrak{su}(V_\rho) \cong \ImH$, of real
        rank exactly~$3$.
        A transversal symmetry generating a 4th independent admissible direction would contradict
        this bound.
        The contradiction shows that $g$ cannot act transversally.
        Therefore $g$ acts vertically.
        Thus, no admissible symmetry can increase the dimension of $\operatorname{Im}\Pi \cap \Ntrl$.
      \end{proof}

    \subsection{Case analysis}
      \label{sec:cases}

      The proof of Lemma~\ref{lem:verticality} proceeds by contradiction and applies uniformly to
      all $g \in \GBI$.
      For clarity, we spell out how the argument applies to the three natural candidate classes of
      additional symmetries discussed in O18 (Problem~2.8 and Remark~2.9) and O23.

      \paragraph{Case 1: Parity $\chi \mapsto -\chi$.}
        This is the response-family covariance established in O18~\cite{Beau2026a22}
        (Proposition~2.6 there).
        Under the hypothesis that the parity orbit $\{\chi, -\chi\}$ lies in a fibre of $\Pi$
        (the open identification of O18 Problem~2.8), it maps that fibre to itself and is
        vertical by definition.
        It does not generate any new direction in $\operatorname{Im}\Pi$.

      \paragraph{Case 2: Phase rotation $\chi \mapsto e^{i\theta}\chi$.}
        The substrate field $\chi$ is real by construction in the Cosmochrony
        framework~\cite{Cosmochrony}.
        A U(1) phase rotation is therefore not a symmetry of $\SBI[\chi]$ for $\theta \notin \pi\mathbb{Z}$,
        and reduces to parity for $\theta = \pi$.
        This class of symmetries does not arise.
        Even if the field were complexified, such a symmetry would act within fibres unless it generated
        new independent admissible neutral traceless directions; but any such directions would be
        excluded by Lemma~\ref{lem:verticality} via the rank-three bound of O23.

      \paragraph{Case 3: Topological symmetries.}
        The Topological Invariants paper~\cite{Beau2026k} establishes that admissible configurations
        are classified by discrete topological invariants (winding number, spin structure) that
        distinguish configurations with distinct invariant values.
        Configurations with different topological invariants are therefore distinguishable by $\Pi$;
        they do not lie in the same fibre and cannot be related by a fibre-preserving symmetry.
        Any topological symmetry that would identify configurations of different winding number would
        contradict the discrete classification of~\cite{Beau2026k}.
        Such symmetries are therefore absent from $\GBI$.

        In each case, the conclusion is consistent with Lemma~\ref{lem:verticality}: no element of
        $\GBI$ acts transversally.

  \subsection{Remark on the non-classificatory character of the argument}
    \label{sec:nonclassif}

    The proof of Lemma~\ref{lem:verticality} does not enumerate all possible symmetries of
    $\SBI$.
    It instead uses the rank-three bound of O23 (conditional on the supplied carrier) as a
    structural obstruction: any transversal symmetry
    would produce an admissible direction of rank greater than~$3$, contradicting that bound.
    This argument is therefore independent of a global classification of $\mathcal{C}_{\mathrm{eff}}$
    and avoids the topological difficulty identified in Remark~4.2 of~\cite{Beau2026k}.

  \section{Main Results}
  \label{sec:main}

  \subsection{Observable rank stability}
    \label{sec:rankstability}

    \begin{theorem}[Observable rank stability]
      \label{thm:rankstability}
      Under the Born--Infeld constraint, the Weil realisation on
      $\mathrm{Heis}_3(\mathbb{Z}/q\mathbb{Z})$, and the supplied spinor carrier
      $V_\rho \cong \mathbb{C}^2$ of O23, the real rank of the admissible neutral traceless
      observable sector $\operatorname{Im}\Pi \cap \Ntrl$ is fixed to~$3$ as a structural invariant of
      the admissible sector, independently of the cardinality of the fibres of~$\Pi$.
      Consequently, the supplied selection rule $\Sigma_c(n_3) = 3$ is insensitive to the fibre
      structure of~$\Pi$, though it remains conditional on the carrier and on the open
      identification of $\Sigma_c$ with the carrier dimension.
      In particular, no admissible enlargement of $\ker\Pi$ can induce any increase in the
      observable rank.
    \end{theorem}

    \begin{proof}
      By Lemma~\ref{lem:decoupling}, the rank of $\operatorname{Im}\Pi \cap \Ntrl$ is independent of
      $|\ker\Pi|$.
      By Lemma~\ref{lem:verticality}, every $g \in \GBI$ acts vertically and therefore does not
      increase this rank.
      By Theorem~\ref{thm:o23} (O23, conditional on the supplied carrier), the rank is fixed
      to~$3$.
      The combination of these three results gives the claim.
      The reading of the threshold $\Sigma_c(n_3) = 3$ as the observable signature of this rank
      is the supplied selection rule of O23~\cite{Beau2026a27}, whose observable identification
      remains open there.
    \end{proof}

  \subsection{Fibre-structure closure and transfer boundary}
    \label{sec:transfer}

    \begin{corollary}[Fibre-structure closure for $\cchi \to \dpair$]
      \label{cor:transfer}
      The observable-rank condition supporting the fibre-side segment
      $\cchi \to \dpair$ is independent of the cardinality of the
      fibres of $\Pi$.
      The absence-of-further-symmetries hypothesis --- parity as the only symmetry of $\SBI$,
      a clause of the open fibre-identification problem (O18 Problem~2.8, Remark~2.9) --- is
      not needed for this fibre-side result:
      it holds for any $\GBI$-symmetric non-injective projection,
      provided the verticality condition of Lemma~\ref{lem:verticality} is satisfied.
      The verticality condition follows from the Born--Infeld admissibility constraint
      and the rank-three bound of O23; it therefore remains conditional on the supplied
      spinor carrier of O23, and the fibre identification itself stays open.
    \end{corollary}

    \begin{proof}
      The fibre-side conditions were formulated in
      O21~\cite{Beau2026a25} Theorem~5.1, under the shell-alignment conjecture and the condition
      that $\Sigma_c(n_3) = 3$.
      O22~\cite{Beau2026a26} proved the shell-alignment conjecture (projection locking forces
      saturation onto a BFS shell).
      Theorem~\ref{thm:rankstability} establishes that $\Sigma_c(n_3) = 3$ is insensitive to
      the fibre structure, conditional on the supplied carrier of O23.
      These two conditions are therefore satisfied independently of fibre
      structure, under the same supplied hypotheses, proving the stated rank stability.
      They do not supply O21's independent cross-substrate growth condition, and hence do not
      establish a native $\dpair \to \bstar$ transfer~\cite{Beau2026sgn}.
    \end{proof}

    \begin{remark}
      Corollary~\ref{cor:transfer} does not remove the numerical open direction.
      The value $\dpair \approx 7.44$ and the shell distance $\mathrm{dist}(n_3^{\mathrm{obs}},
      \mathbb{Z}) \approx 0.20$ at $q = 61$ are numerical results requiring confirmation at larger
      primes $q \in \{101, 151, 211\}$ across multiple conjugate pairs.
      The value $7.44$ uses the production $\log(n+1)$ convention.
      On the same finite windows, translation to $\log n$ gives approximately $6.6$--$6.9$.
      The reciprocal value remains a phenomenological coincidence check, not a structural
      derivation~\cite{Beau2026sgn}.
    \end{remark}

  \section{Discussion}
  \label{sec:discussion}

  \subsection{What O24 closes and what it does not}
    \label{sec:closes}

    The present paper reduces the fibre-structure conditionality in the
    $\cchi \to \dpair$ segment to the supplied carrier and fibre hypotheses.
    The dispositions of the fibre-side problems identified across O21--O23 are:

    \begin{center}
      \begin{tabular}{lll}
        \hline
        Condition & Origin & Status after O24 \\
        \hline
        Shell-alignment conjecture & O21 Conj.~4.1 & Resolved by O22 \\
        $\Sigma_c(n_3) = 3$ & O21 open problem~2 & Supplied selection rule; \\
        & & conditional in O23 (Thm~3.1) \\
        Parity = only symmetry of $\SBI$ & O18 Prob.~2.8, Rem.~2.9 & Replaced by verticality, \\
        & & proved in O24 \\
        Native pair-capacity growth carrier & SGN & Not supplied by O22--O24 \\
        \hline
      \end{tabular}
    \end{center}

    The capacity-to-rate step remains open at the level of new model construction and refuted as
    a native composite for the corpus-defined Heisenberg objects~\cite{Beau2026sgn}.

  \subsection{Relation to O18: necessary versus sufficient}
    \label{sec:o18relation}

    The absence-of-further-symmetries clause --- parity as the only global symmetry of $\SBI$,
    a clause of the open fibre-identification problem (O18 Problem~2.8, Remark~2.9) --- is a
    sufficient condition for the programme.
    It would guarantee that parity fibres are minimal, which implies in particular that no
    additional observable directions are created.

    The present paper identifies the \emph{necessary} condition: verticality of all admissible
    symmetries.
    Verticality is strictly weaker than fibre minimality: a vertical symmetry may enlarge the
    fibre without changing the image.
    The replacement of the sufficient clause by the necessary and sufficient condition
    (O24) therefore strictly extends the scope of the fibre-side result: observable-rank stability
    holds not only when fibres are minimal, but whenever additional fibre structure acts vertically.

  \subsection{Vertical versus transversal non-injectivity}
    \label{sec:vertrans}

    The distinction between vertical and transversal non-injectivity is the structural principle
    underlying the present paper.
    It can be stated as a single invariance property:

    \begin{quote}
      \textbf{Observable rank rigidity under admissible non-injectivity.}
      \emph{The real rank of $\operatorname{Im}\Pi \cap \Ntrl$ is a protected invariant:
        $\ker\Pi$ may change freely; $\operatorname{Im}\Pi \cap \Ntrl$ is rigid.}
    \end{quote}

    \emph{Vertical non-injectivity} means that multiple substrate configurations project to the
    same observable.
    This increases the microscopic degeneracy --- the number of configurations in each fibre ---
    but leaves the observable sector unchanged.
    Its physical effect is an increase in the amplitude of projection residuals, manifest as
    fluctuations in the effective description.

    \emph{Transversal non-injectivity} would mean that $\Pi$ fails to distinguish configurations
    that differ in the observable sector.
    This would collapse independent observable directions and reduce the rank of
    $\operatorname{Im}\Pi \cap \Ntrl$ below~$3$.
    Lemma~\ref{lem:verticality} shows that this case does not arise for any $g \in \GBI$.

    The programme is therefore stable under arbitrary enlargement of the fibres, provided that
    enlargement is vertical.

  \subsection{Fluctuations as kernel-level residuals}
    \label{sec:fluctuations}

    Section~\ref{sec:residuals} identified projection residuals as the physical consequence of
    vertical non-injectivity.
    A larger fibre produces a larger space of substrate configurations that are invisible to
    admissible observables; their contribution appears as fluctuations in the effective
    description.

    Crucially, these fluctuations inhabit the same effective space as the observables.
    They do not create a new sector; they add noise to the existing one.
    The pair-capacity exponent $\dpair$, the saturation rank $n_3$, and the threshold condition
    $\Sigma_c(n_3) = 3$ are all properties of the observable sector and are unaffected by the
    amplitude of projection noise.

    This is consistent with the interpretation developed in O14~\cite{Beau2026a18} Section~6.4,
    where the residual correction $\epsilon(q)$ is proposed as the observable imprint of
    projection residuals, controlled by the phase variance $\mathrm{Var}(\theta)(q)$.
    The present result provides the structural justification for this interpretation: projection
    rank stability is compatible with that interpretation, but does not derive the estimator-level
    correction or its insertion into a reciprocal law.

  \section{Conclusion}
  \label{sec:conclusion}

  The present paper establishes that the observable-rank condition in the fibre-side
  $\cchi \to \dpair$ segment is independent of the cardinality of the fibres of the
  non-injective projection $\Pi$.
  The key result is the Verticality Lemma (Lemma~\ref{lem:verticality}): every Born--Infeld-admissible
  symmetry acts vertically with respect to $\Pi$, so that no admissible symmetry can enlarge
  the real rank of the observable sector beyond~$3$.

  The argument is non-classificatory.
  It does not enumerate all symmetries of $\SBI$ or classify the full space
  $\mathcal{C}_{\mathrm{eff}}$.
  It uses a single structural obstruction: the rank-three bound of
  O23~\cite{Beau2026a27} (Theorem~3.1, conditional on the supplied spinor carrier), which
  ensures that no admissible direction of rank greater than~$3$ can exist within the carrier.

  The conceptual reorientation relative to O18 is:
  \begin{quote}
    \emph{The absence-of-further-symmetries clause would require fibres to be small.
    The present result shows that fibre size is physically irrelevant:
    non-injectivity affects degeneracy, not structure.}
  \end{quote}

  O22, O23, and the present paper reduce the named fibre-structure conditionalities to the
  supplied carrier and fibre hypotheses; the carrier selection, the $\Sigma_c$ identification,
  and the fibre identification of O18 Problem~2.8 remain open.
  They do not derive the separate Heisenberg capacity-to-rate bridge.

  \paragraph{Status update (O25 extension to $q = 601$).}
    That campaign has since been carried out in O25~\cite{Beau2026a29}.
    An extended pair-capacity campaign up to $q = 601$ finds that the reciprocal parameter
    computed from the pair observable remains in the reference $0.11$--$0.13$ range, while the
    underlying \emph{raw} exponent $\dpair(q)$ itself descends into the admissible window
    $[7.4, 10.6]$ as $q$ grows --- reaching $\approx 7.6$ at $q = 601$ --- without requiring the
    finite-size normalization correction used on the earliest (small-$q$) datasets.
    This is a phenomenological consistency check under the cross-substrate prescription, not a
    confirmation of a native transfer.
    Moreover, the production $\log(n+1)$ convention inflates the finite-window exponent relative
    to $\log n$; the asymptotic value $\delta_\infty$ remains open~\cite{Beau2026sgn}.

  \appendix

  \section{Classification of Symmetry Actions: Vertical and Transversal}
  \label{app:examples}

  This appendix collects the three natural candidate classes of additional symmetries of $\SBI$
  and verifies in each case that the action is vertical, supplementing the case analysis of
  Section~\ref{sec:cases}.

  \paragraph{A.1 Parity $\chi \mapsto -\chi$.}
    The Born--Infeld action is even in $\chi$: $\SBI[-\chi] = \SBI[\chi]$.
    Parity therefore belongs to $\GBI$.
    Under the hypothesis that the parity orbit lies in a fibre of $\Pi$ (the open
    fibre-identification problem of O18~\cite{Beau2026a22}, Problem~2.8; evenness alone does
    not force it, O18 Proposition~2.4), $\Pi(-\chi) = \Pi(\chi) = o$.
    Parity then acts vertically.

  \paragraph{A.2 Phase rotation U(1).}
    The substrate field $\chi\colon \mathrm{Heis}_3(\mathbb{Z}/q\mathbb{Z}) \to \mathbb{R}$ is
    real-valued by construction~\cite{Cosmochrony}.
    The Born--Infeld kinetic term $F^2 = (D\chi)^2$ is manifestly real and does not admit a
    U(1) phase symmetry $\chi \mapsto e^{i\theta}\chi$ for $\theta \notin \pi\mathbb{Z}$.
    This class of symmetries is absent from $\GBI$ and requires no further analysis.

  \paragraph{A.3 Topological winding symmetries.}
    The Topological Invariants paper~\cite{Beau2026k} establishes that admissible configurations
    carry discrete topological invariants: charge as a winding invariant (Section~2 of~\cite{Beau2026k}) and spin as
    a covering invariant (Section~3 of~\cite{Beau2026k}).
    Two configurations with distinct winding numbers are distinguished by $\Pi$: they project to
    distinct observables.
    A symmetry that would identify them --- mapping a configuration of winding number $w$ to one
    of winding number $w' \neq w$ --- would contradict the discrete classification of~\cite{Beau2026k} and would not
    be compatible with Born--Infeld admissibility.
    Such symmetries are absent from $\GBI$.

    In all three cases, no transversal action on the admissible observable sector arises.
    The verticality of $\GBI$ established in Lemma~\ref{lem:verticality} is therefore consistent
    with the full catalogue of natural candidate symmetries.

    \clearpage
    \bibliographystyle{plain}
      \bibliography{cosmochrony-bibliography}

\end{document}
